%% ma: L12-B07-0033
%% lop: 12
%% chuong: 2
%% bai: 7
%% dang: 12.7.10
%% loai: TN4
%% muc: TH
%% dap_an: "A"
%% nguon: "CĐ2.6.pdf giáo viên gửi 27/10/2026 (M07, Lê Minh Khôi – Nguyễn Anh Khôi, Chương II Chủ đề 6 Ứng dụng vào các bài toán thực tế), Câu 27"
%% kien_thuc: "Bài 7 (tọa độ điểm khi gắn hệ trục), Bài 8 (trung điểm)"
%% trang_thai: chinh-thuc
%% ghi_chu: "Giáo viên đã duyệt 27/10/2026. Đã bổ sung giả thiết điểm D có hoành độ dương (đề gốc chưa xác định chiều trục Ox). Đề gốc dùng chung giả thiết và hình cho Câu 27–28; đã chép vào từng câu và bổ sung 'D có hoành độ dương' (đề gốc chưa xác định chiều trục Ox)"
\begin{ex}
Để chuẩn bị cho một buổi triển lãm quốc tế, các trang sức có giá trị lớn được đặt bảo mật trong các khối chóp tứ giác đều $S.ABCD$ và đặt lên phía trên một trụ hình hộp chữ nhật $ABCD.A'B'C'D'$ có đáy là hình vuông. Chọn hệ trục tọa độ $Oxyz$ (đơn vị trên mỗi trục là mét) sao cho $A'(0;0;0)$, $A(0;0;1)$, $B(0;0{,}5;1)$ và điểm $D$ có hoành độ dương. Biết rằng ban tổ chức sự kiện dự định dùng các tấm kính cường lực hình tam giác cân có cạnh bên là $60$ cm để lắp ráp thành khối chóp nói trên (hình minh họa).
\begin{center}
\begin{tikzpicture}[scale=3.4,x={(-.55cm,-.45cm)},y={(1cm,0cm)},z={(0cm,1cm)}]
  \path (0,0,0) coordinate (A') (.5,0,0) coordinate (D') (.5,.5,0) coordinate (C') (0,.5,0) coordinate (B')
        (0,0,1) coordinate (A) (.5,0,1) coordinate (D) (.5,.5,1) coordinate (C) (0,.5,1) coordinate (B) (.25,.25,1) coordinate (I) (.25,.25,1.485) coordinate (S);
  \draw (D')--(C')--(B')--(B)--(S)--(D)--(D')  (D)--(C)--(B) (C)--(C') (S)--(C);
  \draw[khuat] (D')--(A')--(B') (A')--(A) (D)--(A)--(B) (S)--(A) (S)--(I) (A)--(C) (B)--(D);
  \path (S) node[above]{$S$} (A) node[above right]{$A$} (B) node[right]{$B$} (C) node[below right]{$C$} (D) node[left]{$D$} (I) node[below right,inner sep=1pt]{$I$}
        (A') node[above right]{$A'$} (B') node[right]{$B'$} (C') node[below]{$C'$} (D') node[left]{$D'$};
\end{tikzpicture}
\end{center}
Gọi $I$ là tâm của hình vuông $ABCD$. Tọa độ của điểm $I$ là
\choice{\True $I(0{,}25;0{,}25;1)$}{$I(0{,}5;0{,}5;1)$}{$I(0;0{,}25;1)$}{$I(0{,}25;0;1)$}
\loigiai{Cạnh hình vuông $AB=0{,}5$; $D(0{,}5;0;1)$, $C(0{,}5;0{,}5;1)$. $I$ là trung điểm $AC$: $I(0{,}25;0{,}25;1)$.}
\end{ex}
